\chapter{Automorphic Forms as Functions on the Adele Group of \texorpdfstring{$\GL(2)$}{GL(2)}}

\label{ch:3}

In Section~\ref{ch:5} we shall realize classical cusp forms not as functions on $\SL(2,\R)$ but as irreducible representations of the adele group of $\GL(2)$ over $\Q$. To make this correspondence explicit we shall (in Section~\ref{ch:4}) collect several facts from the representation theory of $\GL(2)$ over local and global fields. This background is necessary if one wants to understand the precise sense in which the correspondence between cusp forms and representations fails to be one-to-one. It also provides an appropriate setting for the work of Jacquet and Langlands where the classical theory of automorphic forms is attacked from the point of view of group representations.

Our goal in this section is to describe how classical cusp forms correspond first to special functions on the adele group of $\GL(2)$.

\section{Basic Notions}\label{sec:3a}

Throughout this subsection, $p$ will denote an arbitrary place of $\Q$. By $\Q_p$ we shall denote the completion of $\Q$ at $p$ and by $\bG_p$ the local group $\GL(2,\Q_p)$. If $p$ is infinite then $\bG_p$ is $\GL(2,\R)$. If $p$ is finite we let $O_p$ denote the ring of integers of $\Q_p$ and set $\bK_p$ equal to $\GL(2,O_p)$. Each $\bK_p$ is then a maximal compact open subgroup of $\bG_p$. (If $p = \infty$ we take $\bK_p$ to be $\operatorname{O}(2,\R)$.)

The adele ring of $\Q$ will be denoted $\A$ (or $\A(\Q)$ if this reference to the global field is necessary). The group

\[
\bG_A = \GL(2,\A)
\]

is then the direct product of the $\bG_p$ above restricted with respect to $\bK_p$.

Its center is

\[
Z_A = \left\{\begin{pmatrix} t & 0 \\ 0 & t \end{pmatrix}: t \in \A^{\times}\right\}.
\]

By $\bG_Q$ we denote the discrete subgroup $\GL(2,\Q)$ of $\bG_A$. Since $\bG_A$ is unimodular, the quotient space $\bG_Q\backslash\bG_A$ inherits a natural $\bG_A$-invariant measure. The center of $\bG_Q$ is $Z_Q = \left\{\begin{pmatrix} a & 0 \\ 0 & a \end{pmatrix}: a \in \Q^{\times}\right\}$ so $Z_Q\backslash Z_A$ is naturally isomorphic to the idele class group $\Q^{\times}\backslash\A^{\times}$. The measure of the homogeneous space $Z_A\bG_Q\backslash\bG_A$ is finite.

We recall now the principle of strong approximation for $\SL(2)$ over $\Q$. This principle implies (see ``Notes and References'') that

\begin{equation}\label{eq:3.1}
\bG_A = \bG_Q\bG_\infty^+\bK_0',\tag{3.1}
\end{equation}

if $\bG_\infty^+ = \GL^+(2,\R)$, $\bK_0' = \prod_{p<\infty}\bK_p'$, and $\bK_p'$ is any choice of open subgroup of $\bK_p$ such that $\bK_p' = \bK_p$ for almost every $p$ and such that the determinant map from $\bK_p$ to $O_p^{\times}$ is surjective for every $p$. In particular, \eqref{eq:3.1} will hold when

\begin{equation}\label{eq:3.2}
\bK_p^N = \left\{\begin{pmatrix} a & b \\ c & d \end{pmatrix} \in \bK_p: c \equiv 0 \ (\mathrm{mod}\ N)\right\},\tag{3.2}
\end{equation}

$N$ any positive integer. In this case we write $\bK_0^N$ for $\bK_0'$.

The analogue of \eqref{eq:3.1} for $\GL(1)$ is

\begin{equation}\label{eq:3.3}
\A^{\times} = \Q^{\times}\R_+^{\times}\prod_{p<\infty}O_p^{\times}.\tag{3.3}
\end{equation}

This expression is equivalent to the fact that $\Z$ has class number 1.

By a grossencharacter of $\Q$ we shall understand a unitary character of the idele class group of $\Q$, i.e.\ a character of $\A^{\times}$ trivial on $\Q^{\times}$. Observe that by \eqref{eq:3.3} each character of $(\Z/N\Z)^{\times}$ canonically determines a grossencharacter of $\Q$. Indeed suppose $\psi$ is such a character. Then $\psi$ determines a character $\psi_p$ of $O_p^{\times}$ by composition with the natural homomorphism from $O_p^{\times}$ to $(\Z/N\Z)^{\times}$. The product $\prod_{p<\infty}\psi_p$ then describes a character of $\prod_{p<\infty}O_p^{\times}$, hence a character of $\A^{\times}$ (trivial on $\Q^{\times}\R_+^{\times}$) by \eqref{eq:3.3}.

Now fix $f \in S_k(N,\psi)$. Using \eqref{eq:3.1}, and the fact that

\[
\begin{pmatrix} a & b \\ c & d \end{pmatrix} \mapsto \psi_p(a)
\]

determines a character of $\bK_p^N$, we can define a function $\phi_f(g)$ on $\bG_A$ by

\begin{equation}\label{eq:3.4}
\phi_f(g) = f(g_\infty(i))\,j(g_\infty,i)^{-k}\,\psi(k_0)\tag{3.4}
\end{equation}

if $g = \gamma g_\infty k_0$ and $\prod \psi_p = \psi$. This function is well defined since

\begin{equation}\label{eq:3.5}
\bG_Q \cap \bG_\infty^+ \prod \bK_p^N = \Gamma_0(N).\tag{3.5}
\end{equation}

Indeed $f$ is assumed to belong to $S_k(N,\psi)$ and consequently $\phi_f(\gamma g) = \phi_f(g)$ for all $\gamma \in \Gamma_0(N)$.

Observe that if $\psi = 1$, and we regard $\phi_f$ as a function of $\SL(2,\R)$ alone, then $\phi_f$ coincides with the $\phi_f$ of Proposition~\ref{prop:2.1}. In general, using strong approximation, it is not difficult to establish the following proposition.

\begin{proposition}\label{prop:3.1}

The map $f(z)$ goes to $\phi_f(g)$ describes an isomorphism between $S_k(N,\psi)$ and the space of functions $\phi$ on $\bG_A$ satisfying the following conditions:

\begin{enumerate}[(i)]

\item $\phi(\gamma g) = \phi(g)$ for all $\gamma \in \bG_Q$;

\item $\phi(gk_0) = \phi(g)\psi(k_0)$ for all $k_0 \in \prod \bK_p^N$;

\item $\phi(gr(\theta)) = e^{-ik\theta}\phi(g)$ if $r(\theta) = \begin{pmatrix} \cos\theta & -\sin\theta \\ \sin\theta & \cos\theta \end{pmatrix}$;

\item the function $\phi$, viewed as a function of $\bG_\infty^+$ alone, satisfies the differential equation

\begin{equation}\label{eq:3.6}
\Delta\phi = -\frac{k}{2}\left(\frac{k}{2}-1\right)\phi;\tag{3.6}
\end{equation}

\item $\phi(zg) = \phi(gz) = \psi(z)\phi(g)$ for all $z \in Z_A$;

\item $\phi$ is slowly increasing, i.e.\ for every $c > 0$, and compact subset $\omega$ of $\bG_A$, there exist constants $C$ and $N$ such that

\[
\phi\left(\begin{pmatrix} a & 0 \\ 0 & 1 \end{pmatrix}g\right) \leq C|a|^N
\]

for all $g \in \omega$, and $a \in \A^{\times}$ with $|a| > c$: This is the natural adelic version of the classical notion described earlier in the context of Definition~\ref{def:2.2};

\item $\phi$ is cuspidal, i.e.

\[
\int_{\Q\backslash\A}\phi\left(\begin{pmatrix} 1 & x \\ 0 & 1 \end{pmatrix}g\right)\dd x = 0 \quad \text{for almost every } g.
\]

\end{enumerate}

\end{proposition}

We remark that the proof of Proposition~\ref{prop:3.1} involves showing that the ``cuspidal'' (resp. slowly increasing) condition for $\phi_f$ implies that $f(z)$ vanishes (resp. is regular) at every cusp of $\Gamma$.

\begin{definition}\label{def:3.2}

Let $\bL^2(\bG_Q\backslash\bG_A,\psi)$ denote the Hilbert space of measurable functions $\phi$ on $\bG_A$ such that

\begin{enumerate}[(i)]

\item $\phi(\gamma g) = \phi(g)$ for all $\gamma \in \bG_Q$;

\item $\phi(gz) = \phi(g)\psi(z)$ for all $z \in Z_A$; and

\item

\begin{equation}\label{eq:3.7}
\int_{Z_A\bG_Q\backslash\bG_A}|\phi(g)|^2\dd g < \infty.\tag{3.7}
\end{equation}

\end{enumerate}

\end{definition}

It is easy to check that $\phi_f \in \bL^2(\bG_Q\backslash\bG_A,\psi)$ if $f \in S_k(N,\psi)$. Indeed Conditions 3.6(i)-(iii) and (v) imply that $|\phi_f(g)|^2$ is actually a function on $Z_A\bG_Q\backslash\bG_A/\bK_0^N\bK_\infty$. So since

\begin{equation}\label{eq:3.8}
Z_A\bG_Q\backslash\bG_A/\bK_0^N\bK_\infty \cong \Gamma_0(N)\backslash\SL(2,\R)/\operatorname{SO}(2)\tag{3.8}
\end{equation}

(by \eqref{eq:3.1} and \eqref{eq:3.5}) we have

\[
\int_{Z_A\bG_Q\backslash\bG_A}|\phi_f(g)|^2\dd g = \int_{\Gamma_0(N)\backslash\{\operatorname{Im}(z) > 0\}}|f(z)|^2 y^k\frac{\dd x\dd y}{y^2} < \infty,
\]

which proves our claim.

\begin{definition}\label{def:3.3}

An automorphic form on $\GL(2)$ is any function $\phi$ on $\bG_A$ satisfying the following conditions:

\begin{enumerate}[(i)]

\item $\phi(\gamma g) = \phi(g)$ for all $\gamma \in \bG_Q$;

\item for some grossencharacter $\psi$,

\[
\phi(gz) = \phi(zg) = \psi(z)\phi(g) \quad \text{for all } z \in Z_A;
\]

\item

\begin{equation}\label{eq:3.9}
\phi \text{ is right } \bK = \bK_\infty\prod_{p<\infty}\bK_p\text{-finite;}\tag{3.9}
\end{equation}

\item as a function of $\bG_\infty$ alone, $\phi$ is smooth and $\mathfrak{Z}$-finite (here $\mathfrak{Z}$ denotes the center of the universal enveloping algebra of $\bG_\infty$); and

\item $\phi$ is slowly increasing (in the sense of 3.6(vi)).

\end{enumerate}

\end{definition}

If, in addition to (i)-(v), the function $\phi$ satisfies the cuspidal condition 3.6(vii), namely

\[
\int_{\Q\backslash\A}\phi\left(\begin{pmatrix} 1 & x \\ 0 & 1 \end{pmatrix}g\right)\dd x = 0 \quad \text{for a.e.\ } g,
\]

then $\phi$ is a cusp form on $\GL(2)$ (or a $\psi$-cusp form, by 3.9(ii)). The space of all such cusp forms is denoted $A_0(\psi)$.

We have already remarked that the classical holomorphic cusp forms provide special examples of $\psi$-cusp forms in the above sense, and that these cusp forms are ``square integrable on $Z_A\bG_Q\backslash\bG_A$.'' In general, one can show that every $\psi$-cusp form belongs to $\bL^2(\bG_Q\backslash\bG_A,\psi)$. (This statement, of course, is false for automorphic forms which are not cusp forms.)

\begin{remark}\label{rem:3.4}

The subspace of $\bL^2(\bG_Q\backslash\bG_A,\psi)$ consisting of cuspidal functions shall be denoted $\bL_0^2(\bG_Q\backslash\bG_A,\psi)$ and referred to loosely as the space of $\psi$-cusp forms. This terminology is apt because $A_0(\psi)$ coincides with the dense subspace of $\bL_0^2(\bG_Q\backslash\bG_A,\psi)$ consisting of K-finite $\mathfrak{Z}$-finite functions.

\end{remark}

\begin{remark}\label{rem:3.5}

(Concerning the Fourier expansion of an arbitrary $\psi$-cusp form.) The Fourier expansion of a classical cusp form at a cusp $s$ was introduced in Remark~\ref{rem:1.3}. The adelic formulation of this notion is quite simple and extends to arbitrary $\psi$-cusp forms $\phi$ (functions in $\bL_0^2(\bG_Q\backslash\bG_A,\psi)$) as follows.

For almost every $g$ in $\bG_A$,

\[
x \mapsto \phi\left(\begin{pmatrix} 1 & x \\ 0 & 1 \end{pmatrix}g\right)
\]

describes a square-integrable function of the compact group $\Q\backslash\A$. Recall that $\A$ is a self-dual group of all of whose characters are of the form $\tau(bx)$, with $b$ in $\A$,

\[
\tau(x) = \prod_{p<\infty}\tau_p(x_p),
\]

$\tau_\infty(x_\infty) = e^{2\pi ix_\infty}$, and $\tau_p(x_p) = 1$ if and only if $x_p \in O_p$. The characters of $\Q\backslash\A$ are of the form $\tau_\xi(x) = \tau(\xi x)$, where $\xi$ is arbitrary in $\Q$. Consequently we have the Fourier expansion

\begin{equation}\label{eq:3.10}
\phi\left(\begin{pmatrix} 1 & x \\ 0 & 1 \end{pmatrix}g\right) = \sum_{\xi\in\Q}\phi_\xi(g)\tau(\xi x)\tag{3.10}
\end{equation}

for each $\phi$. The $\xi$-th Fourier coefficient $\phi_\xi(g)$ is given by

\begin{equation}\label{eq:3.11}
a_\xi(g) = \phi_\xi(g) = \int_{\Q\backslash\A}\phi\left(\begin{pmatrix} 1 & x \\ 0 & 1 \end{pmatrix}g\right)\overline{\tau(\xi x)}\dd x\tag{3.11}
\end{equation}

and clearly depends on $g$. The cuspidal condition 3.6(vii), as expected, says that the zeroth Fourier coefficient $a_0(g)$ is zero for almost every $g$. Also, for cusp forms of type $\phi_f$, $f \in S_k(N,\psi)$, this general Fourier expansion contains the expansions

\begin{equation}\label{eq:3.12}
f\big|_{[\sigma]_k}(z) = \sum_{n=1}^{\infty}a_n e^{2\pi inz/h}\tag{3.12}
\end{equation}

introduced in Section~\ref{ch:1}. The verification of this fact is straightforward but amusing and therefore we include it below.

\end{remark}

\begin{lemma}\label{lem:3.6}

Suppose $f \in S_k(\SL(2,\Z))$ and $\phi = \phi_f$. Then for each $y > 0$,

\begin{equation}\label{eq:3.13}
\phi_\xi\left(\begin{pmatrix} y & 0 \\ 0 & 1 \end{pmatrix}\right) = \begin{cases} a_n e^{-2\pi ny} & \text{if } \xi = n \in \Z \\ 0 & \text{otherwise.} \end{cases}\tag{3.13}
\end{equation}

Consequently,

\[
\phi\left(\begin{pmatrix} 1 & x \\ 0 & 1 \end{pmatrix}g\right) = \sum_{\xi\in\Q}\phi_\xi(g)\tau(\xi x) = \sum_n a_n e^{2\pi inz} = f(z)
\]

if $g = \begin{pmatrix} y & 0 \\ 0 & 1 \end{pmatrix}$, $y > 0$, and $x \in \R$.

\end{lemma}

\begin{proof}

Suppose first that $\xi \notin \Z$. Then for some prime $p$, and some integer $m > 0$, $\xi = ap^{-m}$, with $a$ in $\Q$ relatively prime to $p$. But 3.9(iii) says that $\phi$ is right $\bK_0 = \prod_{p<\infty}\GL(2,O_p)$ invariant. Therefore, since

\[
\begin{pmatrix} 1 & p^{m-1} \\ 0 & 1 \end{pmatrix} = \left(1_2, 1_2, \dots, \begin{pmatrix} 1 & p^{m-1} \\ 0 & 1 \end{pmatrix}, 1_2, \dots\right)
\]

\hfill \textit{p}-th place

belongs to $\bK_0$, it follows that

\[
\begin{aligned}
\phi_\xi\left(\begin{pmatrix} y & 0 \\ 0 & 1 \end{pmatrix}\right) &= \int_{\Q\backslash\A}\phi\left(\begin{pmatrix} 1 & x \\ 0 & 1 \end{pmatrix}\begin{pmatrix} y & 0 \\ 0 & 1 \end{pmatrix}\begin{pmatrix} 1 & p^{m-1} \\ 0 & 1 \end{pmatrix}\right)\overline{\tau(\xi x)}\dd x\\
&= \int_{\Q\backslash\A}\phi\left(\begin{pmatrix} 1 & x+p^{m-1} \\ 0 & 1 \end{pmatrix}\begin{pmatrix} y & 0 \\ 0 & 1 \end{pmatrix}\right)\overline{\tau(\xi x)}\dd x,
\end{aligned}
\]

where $x+p^{m-1}$ denotes the adele $(x,0,\cdots,0,p^{m-1},0,\cdots)$. So making the change of variable $x$ goes to $x-p^{m-1}$, we have

\begin{equation}\label{eq:3.14}
\phi_\xi\left(\begin{pmatrix} y & 0 \\ 0 & 1 \end{pmatrix}\right) = \tau(\xi p^{m-1})\phi_\xi\left(\begin{pmatrix} y & 0 \\ 0 & 1 \end{pmatrix}\right).\tag{3.14}
\end{equation}

But $\tau(\xi p^{m-1}) = \tau(ap^{-1}) = \tau(0,\dots,0,ap^{-1},0,\dots,0) \neq 1$ by our choice of $\tau$. Therefore \eqref{eq:3.14} implies that $\phi_\xi\left(\begin{pmatrix} y & 0 \\ 0 & 1 \end{pmatrix}\right)$ must vanish identically.

On the other hand, suppose $\xi = n \in \Z$. Then

\[
\begin{aligned}
\phi_\xi\left(\begin{pmatrix} y & 0 \\ 0 & 1 \end{pmatrix}\right) &= \int_{\Q\backslash\A}\phi\left(\begin{pmatrix} y & x \\ 0 & 1 \end{pmatrix}\right)\overline{\tau(nx)}\dd x\\
&= \int_{\Z\backslash\R}\phi\left(\begin{pmatrix} y & x \\ 0 & 1 \end{pmatrix}\right)\overline{\tau(nx)}\dd x\\
&= \int_0^1 f(x+iy)e^{-2\pi inx}\dd x = a_n e^{-2\pi ny}
\end{aligned}
\]

and the proof of \eqref{eq:3.13} is complete. The rest of the lemma follows from the fact that $\tau(x) = e^{2\pi inx}$ whenever $x$ is real.

\end{proof}

If $\Gamma' = \Gamma_0(N)$ the relations between \eqref{eq:3.10} and the classical Fourier expansions \eqref{eq:1.6} are established similarly.

\section{Hecke Operators}\label{sec:3b}

We know from A that elements of $S_k(\Gamma)$ correspond to special functions in the Hilbert space $\bL^2(\bG_Q\backslash\bG_A/\bK_0^N)$. Our task now is to reinterpret the classical Hecke operators $T(p)$ in this adelic setting.

If $\Gamma = \Gamma_0(N)$, let $p$ denote any prime not dividing $N$. Then $\bK_p^N = \bK_p = \GL(2,O_p)$. In $\bG_p$, consider the double coset

\[
\bH_p = \bK_p\begin{pmatrix} p & 0 \\ 0 & 1 \end{pmatrix}\bK_p.
\]

(This is the same thing as $\bK_p\begin{pmatrix} 1 & 0 \\ 0 & p \end{pmatrix}\bK_p$ since $w = \begin{pmatrix} 0 & 1 \\ -1 & 0 \end{pmatrix} \in \bK_p$ and $w\begin{pmatrix} p & 0 \\ 0 & 1 \end{pmatrix}w^{-1} = \begin{pmatrix} 1 & 0 \\ 0 & p \end{pmatrix}$.)

If $\phi$ is any function on $\bG_A$ right invariant by $\bK_p$, let $\tilde{T}(p)\phi$ denote the right convolution over $\bG_p$ of $\phi$ with the characteristic function of $\bH_p$. In particular, for each $\phi$ in $\bL^2(\bG_Q\backslash\bG_A/\bK_0^N)$,

\begin{equation}\label{eq:3.15}
\tilde{T}(p)\phi(g) = \int_{\bH_p}\phi(gh)\dd h.\tag{3.15}
\end{equation}

Furthermore, since the elements of $\bG_\infty^+$ and $\bK_p$ commute (when viewed as adeles of $\GL(2)$ in the obvious way) it is easy to verify that each of the properties 3.6(i)-(vii) is preserved by $\tilde{T}(p)$ and, consequently, that the image of $S_k(\Gamma)$ in $\bL_0^2(\bG_Q\backslash\bG_A/\bK_0^N)$ is mapped into itself by $\tilde{T}(p)$.

\begin{lemma}\label{lem:3.7}

If $\phi = \phi_f$, $f \in S_k(\Gamma)$, then

\[
p^{k/2-1}\tilde{T}(p)\phi = \phi_{T(p)f}.
\]

\end{lemma}

\begin{proof}

From the theory of ``elementary divisors'' in $\bG_p$, it follows that

\[
\bH_p = \bigcup_{b=0}^{p-1}\begin{pmatrix} p & -b \\ 0 & 1 \end{pmatrix}\bK_p \cup \begin{pmatrix} 1 & 0 \\ 0 & p \end{pmatrix}\bK_p,
\]

a disjoint union of right $\bK_p$-cosets. On the other hand, $\phi$ is right $\bK_p$-invariant. Consequently, it follows from \eqref{eq:3.15} that

\[
p^{k/2-1}\tilde{T}(p)\phi(g) = p^{k/2-1}\sum_{b=0}^{p-1}\phi\left(g\begin{pmatrix} p & -b \\ 0 & 1 \end{pmatrix}\right) + p^{k/2-1}\phi\left(g\begin{pmatrix} 1 & 0 \\ 0 & p \end{pmatrix}\right).
\]

(The Haar measure on $\bG_p$ being normalized so that $\bK_p$ has mass 1.)

Now recall that $\phi(g) = \phi_f(g) = f(g_\infty(i))j(g_\infty,i)^{-k}$, if $g = \gamma g_\infty k_0$. Observe that

\[
g_\infty k_0\begin{pmatrix} p & -b \\ 0 & 1 \end{pmatrix} = \gamma'\begin{pmatrix} p^{-1} & bp^{-1} \\ 0 & 1 \end{pmatrix}_{\infty}g_\infty k_0',
\]

with

\[
\gamma' = \left(\begin{pmatrix} p & -b \\ 0 & 1 \end{pmatrix}, \begin{pmatrix} p & -b \\ 0 & 1 \end{pmatrix}, \dots, \begin{pmatrix} p & -b \\ 0 & 1 \end{pmatrix}, \dots\right),
\]

and

\[
k_0' = k_0\left(1, \begin{pmatrix} p & -b \\ 0 & 1 \end{pmatrix}^{-1}, \begin{pmatrix} p & -b \\ 0 & 1 \end{pmatrix}^{-1}, \dots, 1, \begin{pmatrix} p & -b \\ 0 & 1 \end{pmatrix}^{-1}, \dots\right).
\]

\hfill \textit{p}-th place

Thus

\[
\phi\left(g\begin{pmatrix} p & -b \\ 0 & 1 \end{pmatrix}\right) = f\left(\begin{pmatrix} p & -b \\ 0 & 1 \end{pmatrix}^{-1}z\right)p^{-k/2} = f\left(\frac{z+b}{p}\right)p^{-k/2},
\]

similarly

\[
\phi\left(g\begin{pmatrix} 1 & 0 \\ 0 & p \end{pmatrix}\right) = f\left(\begin{pmatrix} 1 & 0 \\ 0 & p \end{pmatrix}^{-1}z\right) = f(pz)p^{k/2},
\]

and consequently

\[
\begin{aligned}
p^{k/2-1}\tilde{T}(p)\phi(g) &= p^{-1}\sum_{b=0}^{p-1}f\left(\frac{z+b}{p}\right) + p^{k-1}f(pz)\\
&= p^{k-1}\sum_{\substack{a > 0\\ ad = p}}\sum_{b=0}^{d-1}f\left(\frac{az+b}{d}\right)d^{-k}\\
&= \phi_{T(p)f}(g),
\end{aligned}
\]

as desired.

\end{proof}

\begin{remark}\label{rem:3.8}

(Why $\GL(2)$ and not $\SL(2)$?) There are numerous technical reasons why $\GL(2)$ (and not $\SL(2)$) is taken as the starting point for the modern theory of automorphic forms. On the one hand the existence of a non-trivial center in $\GL(2)$ makes this group the natural habitat for classical cusp forms of weight $k$ and non-trivial conductor $\psi$. On the other there is the obvious fact that $\begin{pmatrix} p & 0 \\ 0 & 1 \end{pmatrix}$ does not belong to $\SL(2)$. This last fact gives rise to the following awkward situation.

Suppose we realize functions in $S_k(\Gamma)$ as special functions on $\SL(2,\A)$. Then for each prime $p$ it is natural to define the operator $\tilde{T}(p)\phi$ as convolution of $\phi$ with the characteristic function of

\[
\SL(2,O_p)\begin{pmatrix} p & 0 \\ 0 & p^{-1} \end{pmatrix}\SL(2,O_p).
\]

The problem is that this operator does not coincide with $T(p)$.

\end{remark}

\section{Arbitrary Base Fields}\label{sec:3c}

Thus far we have regarded $\GL(2)$ as an algebraic group defined over $\Q$ since this base field clearly suffices for the classical theory as described in Section~\ref{ch:1}. However, other classical examples, such as Hilbert's modular forms, strongly suggest that we deal with automorphic forms on $\GL(2)$ defined over arbitrary global fields $\F$.

So suppose first that $\F$ is an arbitrary algebraic number field, finite over $\Q$. Each place of $\F$ is denoted by $v$ and the completion of $\F$ at $v$ by $\F_v$. As usual, if $v$ is finite, $O_v$ denotes the ring integers of $\F_v$, and $O_v^{\times}$ its group of units. The adeles of $\F$ are again denoted by $\A$, but now we observe that

\begin{equation}\label{eq:3.16}
\A^{\times} = \F^{\times}(\F_\infty^0)\left(\prod_{v\text{ finite}}O_v^{\times}\right)\tag{3.16}
\end{equation}

if and only if the class number of $\F$ is one. (Here $\F_\infty^0$ denotes the connected component of the archimedean component of $\F$, the product of the completions $\F_v$ as $v$ ranges over the finite set of infinite places of $\F$.)

Regarding $\GL(2)$ as an algebraic group over $\F$ we write $\bG_A$ for the corresponding adelized group $\GL(2,\A_\F)$. Similarly $\bG_\F = \GL(2,\F)$ and $\bG_v = \GL(2,\F_v)$. We then define a cusp form on $\GL(2)$ over $\F$ by suitably modifying the conditions of Definition~\ref{def:3.3}. In particular, such a form satisfies the following conditions:

\begin{enumerate}[(i)]

\item $\phi(\gamma g) = \phi(g)$ for all $\gamma \in \bG_\F$;

\item for some grossencharacter $\psi$ of $\F$,

\begin{equation}\label{eq:3.17}
\phi(gz) = \psi(z)\phi(g) \quad \text{for all } z \in Z_{\A_\F};\tag{3.17}
\end{equation}

\item $\phi(g)$ is right $\bK = \bK_\infty\prod_{v\text{ finite}}\bK_v$ finite if $\bK_v = \GL(2,O_v)$;

\item viewed as a function of $\bG_\infty = \prod_{v\text{ infinite}}\bG_v$, $\phi$ is smooth and $\mathfrak{Z}_\infty$-finite ($\mathfrak{Z}_\infty$ denotes the center of universal enveloping algebra of $\bG_\infty$);

\item $\phi$ is slowly increasing; and

\item

\[
\int_{\F\backslash\A}\phi\left(\begin{pmatrix} 1 & x \\ 0 & 1 \end{pmatrix}g\right)\dd x = 0 \quad \text{for a.e.\ } g.
\]

\end{enumerate}

As before, the space $A_0(\psi)$ of such forms is characterized as the dense subspace of $\bL_0^2(\bG_\F\backslash\bG_\A,\psi)$ consisting of K-finite $\mathfrak{Z}_\infty$-finite functions.

\begin{remark}\label{rem:3.9}

Suppose $\F$ is a function field in one variable over a finite field of constants. Then the notion of an automorphic cusp form on $\GL(2)$ over $\F$ is particularly simple since in this case there are no infinite places. In fact the space $A_0(\psi)$ is precisely the space of K-finite vectors in $\bL_0^2(\bG_\F\backslash\bG_\A,\psi)$ and conditions (iv) and (v) of \eqref{eq:3.17} may be discarded. (Any K-finite cuspidal function on $\bG_\F\backslash\bG_\A$ satisfying 3.17(ii) is automatically compactly supported (modulo $Z_A$) if $\F$ is a function field.)

\end{remark}

\begin{example}\label{exa:3.10}

(Hilbert modular forms.) Suppose $\F = \Q(\sqrt{d})$ with $d > 0$. For simplicity suppose also that $\F$ has class number 1. Then it can be shown that

\begin{equation}\label{eq:3.18}
\GL(2,\A_\F) = \bG_\F\bG_\infty^0\prod_{v\text{ finite}}\GL(2,O_v)\tag{3.18}
\end{equation}

(using \eqref{eq:3.16} and ``strong approximation'' for $\SL(2)$ over $\F$). Therefore, since $\bG_\infty^0 = \GL^+(2,\R)\times\GL^+(2,\R)$, and (essentially)

\begin{equation}\label{eq:3.19}
\left(\bG_\infty^0\prod\bK_v\right)\cap\bG_\F = \SL(2,O_\F),\tag{3.19}
\end{equation}

where $O_\F = \Z(\sqrt{d})$ denotes the ring of integers of $\F$, we have

\begin{equation}\label{eq:3.20}
Z_A\bG_\F\backslash\bG_\A/\prod\bK_v \cong \SL(2,O_\F)\backslash\SL(2,\R)\times\SL(2,\R).\tag{3.20}
\end{equation}

Since $\SL(2,\R)\times\SL(2,\R)$ acts naturally (and transitively) on $\{\operatorname{Im}(z_1) > 0\}\times\{\operatorname{Im}(z_2) > 0\}$, with isotropy group $\operatorname{SO}(2)\times\operatorname{SO}(2)$, we also have

\[
Z_A\bG_\F\backslash\bG_\A/(\prod\bK_v)\bK_\infty \cong \SL(2,O_\F)\backslash\{\operatorname{Im}(z_1) > 0\}\times\{\operatorname{Im}(z_2) > 0\}.
\]

We recall that a Hilbert modular (cusp) form (of weight $k$) is a holomorphic function in $\{\operatorname{Im}(z_1) > 0\}\times\{\operatorname{Im}(z_2) > 0\}$ which is bounded and such that

\[
f\left(\frac{az_1+b}{cz_1+\overline{d}}, \frac{\overline{a}\,\overline{z}_2+\overline{b}}{\overline{c}\,\overline{z}_2+\overline{d}}\right) = (cz_1+d)^k(\overline{c}\,\overline{z}_2+\overline{d})^k f(z_1,z_2)
\]

whenever $\begin{pmatrix} a & b \\ c & d \end{pmatrix}$ belongs to $\SL(2,O_\F)$. Here $\overline{a}$ denotes the conjugate of $a$ in $\Q(\sqrt{d})$. Using \eqref{eq:3.18} and \eqref{eq:3.19} it is not difficult to show that any such $f(z_1,z_2)$ corresponds to a function $\phi(g)$ on $\GL(2,\A_\F)$ satisfying the following conditions: $\phi(g) \in \bL_0^2(\bG_\F\backslash\bG_\A,1)$, $\phi$ is right $\prod_{v\text{ finite}}\bK_v$-invariant, $\phi$ is an eigenfunction of $\mathfrak{Z}_\infty$, and $\phi(gr(\theta_1)r(\theta_2)) = e^{-ik\theta_1}e^{-ik\theta_2}\phi(g)$ for all $(r(\theta_1),r(\theta_2)) \in \bK_\infty$. Thus $f(z_1,z_2)$ corresponds to a special automorphic cusp form on $\GL(2)$ over $\F$.

\end{example}

\section*{Notes and References}

\addcontentsline{toc}{subsection}{Notes and References}

Suppose we are given an algebraic number field $\F$ and an algebraic group defined over $\F$. If $S$ is some finite set of places and $\bG_S$ is the $S$-component of $\bG_A$ then the problem of strong approximation asks whether $\bG_\F\bG_S$ is dense in $\bG_A$. For $S = S_\infty$ (the set of infinite places of $\F$) and $\bG = \SL(2)$, the answer is yes. This means that for each $g_1$ in $\SL(2,\A(\Q))$ we can find $\gamma_1$ in $\SL(2,\Q)$ such that the $p$-component of $\gamma_1^{-1}g_1$ belongs to $\SL(2,O_p)$ for each $p$. On the other hand,

\[
\A_\F^{\times}/\F^{\times}\F_\infty^0\prod_{v\text{ finite}}O_v^{\times}
\]

is naturally isomorphic to the ideal class group of $\F$ for any number field $\F$. Consequently, taking $\F = \Q$, \eqref{eq:3.3} results. In particular, strong approximation for $\SL(2)$, together with the identity

\[
g = \begin{pmatrix} r & 0 \\ 0 & 1 \end{pmatrix}g_1\begin{pmatrix} u & 0 \\ 0 & 1 \end{pmatrix},
\]

where $g$ in $\GL(2,\A)$ is arbitrary, and $\det g = ru$, with $r \in \Q^{\times}$ and $u \in (\Q_\infty)_+^{\times}\prod_{p<\infty}O_p^{\times}$, indeed implies \eqref{eq:3.1}.

The problem of strong approximation is discussed in \cite{Kneser}. Some general features of automorphic forms on $\GL(2)$ over number fields and function fields are described in \cite{JacquetLanglands}, \cite{GodementJacquet}, and \cite{Weil2}.
